\section[Complete asymptotics for the normalizing constant of the high-dimensional matrix Bingham distribution]{Complete asymptotics for the normalizing constant\\ of the high-dimensional matrix Bingham distribution}
\label{sec_matrix_bingham}

A \textit{partition} $\kappa = (\kappa_1,\ldots,\kappa_d)$ is a vector of integers such that $\kappa_1 \ge \cdots \ge \kappa_d \ge 0$. The non-zero~$\kappa_j$~are called the \textit{parts} of $\kappa$. The \textit{length} of $\kappa$, denoted by $\ell(\kappa)$, is the number of non-zero~$\kappa_j$,~$j=1,\ldots,d$, and the \textit{weight} of $\kappa$ is $|\kappa| = \kappa_1+\cdots+\kappa_d$.

Define the \textit{shifted factorial}, $(a)_k = a(a+1)(a+2) \cdots (a+k-1)$, where $a \in \mathbb{R}$ and ${k=0,1,2,\ldots}$; in particular, $(a)_0 \equiv 1$ since the defining product is empty. For $a \in \mathbb{R}$ and a partition ${\kappa = (\kappa_1,\ldots,\kappa_d)}$, the {\it partitional shifted factorial} is %defined as
\begin{equation}
\label{eq_partit_fact}
(a)_\kappa = \prod_{i=1}^d \left(a-\frac{1}{2}(i-1)\right)_{\kappa_i}.
\end{equation}

Assume that \smash{$\R^{d \times d}_\sym$} denote the space of real symmetric $d \times d$ matrices. For a partition $\kappa = (\kappa_1,\ldots,\kappa_d)$ and \smash{$\Sigma \in \R^{d \times d}_\sym$}, the \textit{zonal polynomial} $C_\kappa(\Sigma)$ is an eigenfunction of the Laplace--Beltrami operator in the eigenvalues of $\Sigma$ (see \cite{James68}, \cite[Chapter 7]{Muirhead} and \cite{Richards10}). Further, $C_\kappa(\Sigma)$~depends only on the eigenvalues of $\Sigma$, and is symmetric in those eigenvalues.

Tables of the polynomials $C_\kappa(\Sigma)$ for $|\kappa| \le 12$ were provided by James \cite[pp.~498--499]{James64} and Parkhurst and James \cite{ParkhurstJames}. Jiu and Koutschan \cite{JiuKoutschan} have provided software to calculate rapidly the zonal polynomials to high values of $|\kappa|$, so the asymptotic expansions obtained in this article can be calculated for practical values of $m$.

Let $X \in V_{d,p}$ be a random matrix having the matrix Bingham distribution with the density function $\phi(x;A,\Sigma)$ as in \eqref{eq_mBingham_pdf}. It is well known that the normalizing constant $\Phi_{d,p}(A,\Sigma)$ can be expressed as a generalized hypergeometric function of two matrix arguments, denoted by~${}_0F_0(A,\Sigma)$, or by the zonal polynomial expansion of that function,
\begin{equation}
\label{eq_normaliz_const_B}
\Phi_{d,p}(A,\Sigma) = {}_0F_0(A,\Sigma) = \sum_{k=0}^\infty \frac{1}{k!} \sum_{|\kappa|=k} \frac{C_\kappa(A) C_\kappa(\Sigma)}{C_\kappa(I_d)},
\end{equation}
where the inner sum is over all partitions $\kappa$ such that $\ell(\kappa) \le p$. This formula for $\Phi_{d,p}(A,\Sigma)$ in terms of ${}_0F_0(A,\Sigma)$ dates to De Waal \cite{DeWaal}; cf.\ Bingham \cite{Bingham}, Chikuse \cite[p.~108]{Chikuse}. By applying the methods of \cite[Theorem~6.3]{GrossRichards}, we find that the series \eqref{eq_normaliz_const_B} converges absolutely for all $(A,\Sigma)$ and uniformly on compact regions.

For the special case in which $A = I_p$, on applying to \eqref{eq_normaliz_const_B} the identity \eqref{eq_kappa_invar}, written in the form $C_\kappa(I_p)/C_\kappa(I_d) = (p/2)_\kappa/(d/2)_\kappa$, we obtain the normalizing constant in terms of a confluent hypergeometric function of matrix argument,
%\label{eq_normaliz_const_hgf}
\[
\Phi_{d,p}(I_p,\Sigma) = \sum_{k=0}^\infty \frac{1}{k!} \sum_{|\kappa|=k} \frac{(p/2)_\kappa}{(d/2)_\kappa} C_\kappa(\Sigma) = {}_1F_1\left(\frac12 p;\frac12 d;\Sigma\right),
\]
as shown by various authors, e.g., Bingham \cite{Bingham}, Chikuse \cite[p.~33]{Chikuse} and Muirhead \cite[p.~288]{Muirhead}.

Similar to \cite[equation~(2.4)]{BagyanRichards}, when $d$ is large, our asymptotic approximation to $\Phi_{d,p}(A,\Sigma)$ is the sum of the terms up to degree $m-1$ of the series \eqref{eq_normaliz_const_B}, i.e.,
\begin{equation}
\label{eq_trunc_expan}
\Phi_{d,p}(A,\Sigma) \approx \sum_{k=0}^{m-1} \frac{1}{k!} \sum_{|\kappa|=k} \frac{C_\kappa(A) C_\kappa(\Sigma)}{C_\kappa(I_d)},
\end{equation}
$m \ge 2$, and now we study the remainder series
\begin{equation}
\label{eq_remainder_term}
\Phi_{d,p;m}(A,\Sigma) = \sum_{k=m}^\infty \frac{1}{k!} \sum_{|\kappa|=k} \frac{C_\kappa(A) C_\kappa(\Sigma)}{C_\kappa(I_d)}.
\end{equation}

In the sequel, we will encounter the constants \smash{$\alpha_p = (2\pi)^{-(p-1)/4p} p^{1/4p}$} and $\gamma_1 = \bigl(\sqrt{3}+1\bigr)/2$, arising in Lemmas \ref{lemma_factorial_ineq} and \ref{lemma_ineqs}, respectively. Also, if $a_1,\ldots,a_p$ are the eigenvalues of $A$, then we~define $A_+ = \diag(|a_1|,\ldots,|a_p|)$, the diagonal matrix with diagonal entries $|a_1|,\ldots,|a_p|$. In~the following result, the proof of which is given in Section \ref{sec_proofs_Phi_bounds}, we extend a result in~\cite{BagyanRichards} by~deriving an upper bound for $\Phi_{d,p;m}(A,\Sigma)$ in terms of the Frobenius norm, \smash{$\|\Sigma\| = \bigl[\tr\bigl(\Sigma^2\bigr)\bigr]^{1/2}$}.

\begin{Theorem}
\label{thm_Phi_m_bound_Bingham}
Let $A \in \R^{p \times p}_\sym$, \smash{$\Sigma \in \R^{d \times d}_\sym$}, and suppose that \smash{$\|\Sigma\| \le \gamma_0 d^{r/2}$}, where $\gamma_0 > 0$ and $0 \le r < 1$. Then for all $m \ge 2$,
\begin{align}
|\Phi_{d,p;m}(A,\Sigma)| \le{}& \alpha_p \sum_{k=m}^\infty \frac{\bigl[\gamma_0 \gamma_1 p^{1/2} (\tr A_+) d^{-(1-r)/2}\bigr]^k}{(k!)^{1/2}} \label{eq_Phi_m_bound} \\
\le{}& \alpha_p (4 {\rm e}/\pi)^{1/4} ({\rm e}/m)^{(m/2)-(1/4)}
\bigl[\gamma_0 \gamma_1 p^{1/2} (\tr A_+) d^{-(1-r)/2}\bigr]^m \nonumber \\
& \times \exp\bigl(c_m \bigl[\gamma_0 \gamma_1 p^{1/2} (\tr A_+) d^{-(1-r)/2}\bigr]^2/2\bigr), \label{eq_Phi_m_boundw}
\end{align}
where $c_m = \bigl(1+m^{-1}\bigr)^{-m} {\rm e}$. In particular, $\Phi_{d,p;m}(A,\Sigma) = O\bigl(d^{-(1-r)m/2}\bigr)$ as $d \to \infty$.
\end{Theorem}
